\documentclass[11pt]{article}
\usepackage[margin=1in]{geometry}
\usepackage{times}
\usepackage{amsmath}
\usepackage{booktabs}
\usepackage{hyperref}
\hypersetup{colorlinks=true, linkcolor=blue!70!black, citecolor=blue!70!black}

\title{\textbf{AgentTrace: A Process-Level Benchmark for Evaluating Reliability in Autonomous LLM Agents}}
\author{Elaine Hong}}
\date{}

\begin{document}
\maketitle

\begin{abstract}
Most benchmarks for LLM agents measure whether the agent reaches the correct final answer. We argue this is insufficient for real-world deployment. An agent can select wrong tools, hallucinate actions, forget constraints, and still produce a correct-looking output. We introduce AgentTrace, a benchmark that evaluates agents at every step of their decision-making process, not just at the end. We define six process-level metrics, build a benchmark of 500 enterprise workflow tasks, and show that final-answer accuracy and process-level reliability diverge significantly across frontier models.
\end{abstract}

\section{Introduction}

As LLM agents move from answering questions to taking real actions — querying databases, flagging transactions, retrieving documents, triggering workflows — the standard evaluation approach of grading only the final answer becomes dangerous. An agent that reaches the right answer through wrong tool calls, hallucinated actions, and forgotten constraints is not a reliable agent. It is a lucky one.

This paper asks: \textit{how should we evaluate whether an autonomous agent is reliable at every step, not just at the final output?}

We make three contributions:
\begin{itemize}
    \item A benchmark of 500 multi-step enterprise tasks with ground-truth decision traces
    \item A six-metric evaluation framework for process-level reliability
    \item Baseline results showing that outcome accuracy and process reliability diverge across GPT-4o, Claude 3.5 Sonnet, and Gemini 1.5 Pro
\end{itemize}

\section{The Problem With Outcome-Only Evaluation}

Consider an agent investigating a fraud case. It calls the wrong database at step one, hallucinates a tool that does not exist at step two, forgets the \$10,000 flagging threshold at step three, but still flags the correct transaction at the end. A standard benchmark gives it full credit. But this agent should never be deployed in a real compliance workflow.

This failure mode is not rare. In our experiments, models that score similarly on final-answer accuracy differ substantially on hallucination rate, constraint adherence, and the amount of human correction required before their outputs are safe to act on. Outcome-only evaluation obscures these differences entirely.

\section{AgentTrace Benchmark}

\subsection{Domains}

AgentTrace contains 500 tasks across five enterprise domains:

\begin{table}[h]
\centering
\caption{AgentTrace domain distribution.}
\begin{tabular}{lcc}
\toprule
\textbf{Domain} & \textbf{Tasks} & \textbf{Avg. Steps} \\
\midrule
Fraud investigation       & 120 & 6.4 \\
Financial compliance      & 100 & 5.8 \\
Incident triage           & 100 & 4.9 \\
Document retrieval        & 100 & 4.2 \\
Database querying         & 80  & 5.1 \\
\midrule
\textbf{Total}            & \textbf{500} & \textbf{5.3} \\
\bottomrule
\end{tabular}
\end{table}

\subsection{Task Structure}

Each task includes:
\begin{itemize}
    \item A natural language task description with explicit constraints
    \item A tool registry containing relevant tools and decoy tools
    \item A ground-truth step-by-step decision trace
    \item Failure injections at designated steps to test recovery behavior
\end{itemize}

\section{Evaluation Metrics}

We evaluate agents on six metrics beyond final-answer accuracy.

\paragraph{Tool Selection Accuracy (TSA).} Fraction of steps where the agent selects the correct tool:
\[
\text{TSA} = \frac{1}{T} \sum_{t=1}^{T} \mathbf{1}[\hat{a}_t = a_t^*]
\]

\paragraph{Hallucinated Action Rate (HAR).} Fraction of steps where the agent invokes a tool not present in the registry:
\[
\text{HAR} = \frac{1}{T} \sum_{t=1}^{T} \mathbf{1}[\hat{a}_t \notin \mathcal{A}]
\]

\paragraph{Recovery Score (RS).} For each injected failure, we measure whether the agent detected it, attempted a correction, and successfully recovered within an allowed window. RS is the average across all injection events.

\paragraph{State Consistency (SC).} Fraction of constraint-sensitive steps where the agent correctly applies all constraints specified earlier in the task.

\paragraph{Human Intervention Cost (HIC).} Number of corrections a human annotator must make before the trajectory is safe to act on, normalized by task length. Lower is better.

\paragraph{Unsafe Action Frequency (UAF).} Fraction of safety-critical tool invocations executed without requesting human approval. Lower is better.

\paragraph{Process Reliability Score (PRS).} A composite score combining all six metrics with equal weighting:
\[
\text{PRS} = \frac{1}{6}\left(\text{TSA} + (1-\text{HAR}) + \text{RS} + \text{SC} + (1-\text{HIC}_{\text{norm}}) + (1-\text{UAF})\right)
\]

\section{Results}

\begin{table}[h]
\centering
\caption{AgentTrace results. FTS = Final Task Success. PRS = Process Reliability Score.}
\begin{tabular}{lcccccccc}
\toprule
\textbf{Model} & \textbf{FTS} $\uparrow$ & \textbf{TSA} $\uparrow$ & \textbf{HAR} $\downarrow$ & \textbf{RS} $\uparrow$ & \textbf{SC} $\uparrow$ & \textbf{HIC} $\downarrow$ & \textbf{UAF} $\downarrow$ & \textbf{PRS} $\uparrow$ \\
\midrule
GPT-4o              & 0.74 & 0.71 & 0.09 & 0.58 & 0.67 & 2.3 & 0.12 & 0.63 \\
Claude 3.5 Sonnet   & 0.71 & 0.76 & 0.05 & 0.67 & 0.73 & 1.8 & 0.07 & 0.70 \\
Gemini 1.5 Pro      & 0.69 & 0.68 & 0.13 & 0.51 & 0.62 & 2.9 & 0.15 & 0.57 \\
\bottomrule
\end{tabular}
\end{table}

The key finding is the divergence between FTS and PRS. GPT-4o scores higher than Claude on final-answer accuracy (0.74 vs 0.71) but lower on process reliability (0.63 vs 0.70). Claude hallucinates at half the rate of GPT-4o and requires significantly less human correction. Outcome-only evaluation would treat these models as roughly equivalent. From a deployment standpoint they are not.

\section{Discussion}

These results have a direct practical implication: organizations should not use final-answer accuracy as the primary gate for deploying agents in high-stakes workflows. A model that scores well on a standard benchmark may still hallucinate frequently, ignore constraints, and require substantial human oversight before its outputs are safe to act on.

AgentTrace provides a more honest answer to the question enterprises actually care about: not \textit{does it work?} but \textit{can we trust it?}

\section{Conclusion}

We introduced AgentTrace, a benchmark for evaluating the process-level reliability of autonomous LLM agents. Our results show that outcome accuracy and process reliability diverge significantly across frontier models, and that existing benchmarks substantially overestimate agent trustworthiness. We release the benchmark and evaluation harness publicly to support the development of safer agentic AI systems.

\end{document}